% 本文件由 LaTeX 排版 Agent 根据有效 translation.json 自动生成。
% author=Justin Martyr; work_id=0129; title=劝勉希腊人书

\chapter{劝勉希腊人书}

% structure: p0001 source=h1 target=omitted reason=page-title-already-rendered-by-parent

% structure: p0002 source=h2 target=section reason=relative-heading-rank; base=h2

\section{第一章 劝勉希腊人的理由}

我开始向你们，希腊人啊，作这一番劝勉之言时，我祈求天主，使我知道该向你们说什么，也使你们摆脱惯常的好辩习性，并从你们祖先的错误中解脱出来，如今能选择有益之事；不要以为，虽然你们先前认为绝对无益的事物如今在你们看来似乎有用，就是冒犯了你们的祖先。因为对事物进行准确的考察，以更深入的探究追问真理，常常显明那些曾被奉为卓越的东西，实际上是另一回事。既然我们打算讨论真宗教（我认为，对于那些希望安然度过此生、因着今生结束后将要到来的审判——这审判不仅由我们按天主的祖先，即先知和立法者所宣告，也由你们中间被视为智慧的人，不仅是诗人，还有哲学家所宣告，他们曾在你们中间声称自己获得了真实而神圣的知识——的人而言，没有什么比这更宝贵），我认为最好首先考察宗教的教师，无论是我们的还是你们的：他们是谁，多么伟大，生活在什么时代；以便那些先前从父辈接受虚假宗教的人，如今觉察到这一点时，能从那种根深蒂固的错误中解脱出来；并且我们能清楚而显明地证明，我们自己所追随的是按天主的祖先的宗教。\par

% structure: p0004 source=h2 target=section reason=relative-heading-rank; base=h2

\section{第二章 诗人不适合作宗教教师}

那么，你们希腊人啊，你们称谁为你们的宗教教师呢？是诗人吗？向了解诗人的人这样说，对你们的事无益；因为他们知道，他们编造了多么可笑的神谱——我们可以从你们最杰出、堪称诗人之王的荷马那里得知。因为他首先说，诸神最初是从水产生的；他这样写道：——\par

\Quote{两者都是诸神的起源——海洋和他们的母亲忒堤斯}\par

然后我们还必须提醒你们，他关于那位你们视为众神之首、并常称他为\Quote{神人之父}者所说的话；因为他说道：——\par

\Quote{宙斯，他是战争的分配者。}\par

事实上，他不仅说他是军队的战争分配者，而且通过他的女儿成为特洛亚人背誓的原因；荷马描绘他陷入情网，痛苦抱怨，自怜自艾，被其他神祇阴谋陷害，有时为他的儿子呼喊：——\par

\Quote{哀哉！我最亲爱的人倒了！\newline{}萨尔珀冬，被帕特罗克洛斯击败，倒了。\newline{}命运如此。}\par

而另一次关于赫克托耳：——\par

\Quote{啊！我看见一位勇士为我所爱，\newline{}在伊利昂城下被驱赶，\newline{}我为赫克托耳悲伤。}\par

至于其他神祇反对宙斯的阴谋，他们知道这些词句：\Quote{当其他奥林匹斯神——朱诺、尼普顿和密涅瓦——想要捆绑他时。}若非有福的神们惧怕那位众神称为布里亚柔斯的神，宙斯早已被他们捆绑。关于他放纵的情欲，我们必须用他原话提醒你们。因为他说宙斯这样对朱诺说：——\par

\Quote{因为从未有女神或妇人，\newline{}向我胸中倾注如此充沛的爱潮；\newline{}我从未这样爱过伊克西翁的妻子，\newline{}也未爱过甜美的阿克里西俄斯的达那厄，\newline{}她生了珀耳修斯，人类最高贵的后裔；\newline{}也未爱过福尼克斯美丽的女儿，她生了\newline{}弥诺斯——无人能及的上界诸神，\newline{}和拉达曼提斯；也未爱过塞墨勒，\newline{}也未爱过阿尔克墨涅，她在忒拜生了\newline{}勇敢的赫拉克勒斯；尽管我的儿子\newline{}由塞墨勒所生是巴克斯，人类的欢乐；\newline{}也未爱过金发的刻瑞斯，或高坐宝座的\newline{}勒托在天上；不——甚至不如\newline{}我现在爱你，我的灵魂感知到\newline{}被强烈渴望的甜蜜所淹没。}\par

现在适宜提及，人们能从荷马的作品中了解到其他神祇的情况，以及他们从人类手中所受的苦。因为他说阿瑞斯和阿佛洛狄忒被狄俄墨得斯所伤，他还叙述了许多其他神祇所受的苦。因为我们可以从狄俄涅安慰她女儿的话中得知；她对她说道：——\par

\Quote{忍耐，最亲爱的孩子；虽被极力逼迫，\newline{}克制你的愤怒：我们这些住在天上的，\newline{}常须忍受凡人的伤害；而我们也\newline{}\newline{}彼此施加许多痛苦：\newline{}阿瑞斯受过苦；被阿洛欧斯的儿子们，\newline{}奥托斯和厄菲阿尔忒斯，牢牢捆绑，\newline{}他在铜镣中躺了十三个月：\newline{}赫拉也受过苦，当安菲特律翁的儿子\newline{}将一支三叉箭射穿她的右乳：\newline{}她所受的苦是可怕而前所未闻的，\newline{}伟大的哈得斯自己也感受到了那支刺痛的箭，\newline{}当那个持埃癸斯的宙斯的同一儿子\newline{}在地狱的大门口袭击他，\newline{}给他造成了最剧烈的苦痛；被痛苦刺穿，\newline{}他呻吟着来到高耸的奥林匹斯，到宙斯的宫殿；\newline{}那支致命的箭仍留在\newline{}他深深的肩头，使他的灵魂悲伤。}\par

但如果提醒你们诸神彼此对战的战斗是合适的，你们自己的诗人自己会叙述，说道：——\par

\Quote{这就是诸神在战斗中相遇的冲击；\newline{}因为那里，向国王尼普顿对面站立的\newline{}是福玻斯·阿波罗，带着他锐利的箭；\newline{}蓝眼的帕拉斯向战神；\newline{}向赫拉，狄安娜，天上的女弓箭手，\newline{}福玻斯的姐妹，金箭女王。\newline{}机敏的赫耳墨斯，助人之神，面对勒托。}\par

这些以及诸如此类的事情，是荷马教导你们的；不仅是荷马，还有赫西俄德。因此，如果你们相信你们最杰出的诗人——他们给出了你们神祇的家谱——，那么你们必须要么认为诸神就是这样的存在，要么相信根本没有神。\par

% structure: p0020 source=h2 target=section reason=relative-heading-rank; base=h2

\section{第三章 泰勒斯学派的意见}

如果你拒绝引用诗人，因为你说他们被允许编造神话，并以神话的方式叙述许多关于神明的、与真理相去甚远的事情，你以为你还有其他宗教导师吗？或者说，你认为他们自己是如何学到你们的这种宗教的？因为任何人都不可能知晓如此伟大而神圣的事理，除非他们自己首先从受启示者那里学到这些。你无疑会说：\Quote{圣贤与哲学家。}因为每当有人引用你们的诗人关于神明的见解时，你习惯逃向他们，如同逃向一座坚固的堡垒。因此，既然我们应该从古人、从最早的那些人开始，我就从那里出发，提出每个人的见解，这些见解比诗人的神学更为可笑。因为\Place{米利都}的\Person{泰勒斯}，在自然哲学研究中领先，宣称水是万物的第一原理；他说万物由水而生，也归于水。在他之后，同样来自\Place{米利都}的\Person{阿那克西曼德}说，无限是万物的第一原理；因为万物确实由此产生，也归于消亡。第三，\Person{阿那克西美尼}——他也来自\Place{米利都}——说气是万物的第一原理；他说万物由气产生，也归于气。来自\Place{梅塔蓬图姆}的\Person{赫拉克利特}和\Person{希帕索}说，火是万物的第一原理；因为万物由火而出，也归结于火。\Place{克拉佐美尼}的\Person{阿那克萨戈拉}说，同质部分是万物的第一原理。\Person{阿凯劳斯}，\Person{阿波罗多罗斯}之子，一位\Place{雅典}人，说无限的气及其密度与稀薄是万物的第一原理。所有这些，从\Person{泰勒斯}开始形成一个系列，遵循他们自称为物理学的哲学。\par

% structure: p0022 source=h2 target=section reason=relative-heading-rank; base=h2

\section{毕达哥拉斯与伊壁鸠鲁的观点}

然后，在另一个起点之后的规律继起中，\Place{萨摩斯}的\Person{毕达哥拉斯}，\Person{穆涅萨库斯}之子，称数及其比例与和谐，以及由两者组成的元素为第一原理；他还包括一与不定之二。\Place{雅典}人\Person{伊壁鸠鲁}，\Person{尼奥克勒斯}之子，说存在物的第一原理是可被理性感知的物体，不容虚空，非受造，不可毁坏，既不能被打碎，也不能有任何部分的形式改变或改变，因此是可被理性感知的。\Place{阿格里真托}的\Person{恩培多克勒}，\Person{梅顿}之子，主张有四种元素——火、气、水、土；以及两种基本力量——爱与恨，前者是结合的力量，后者是分离的力量。那么你看到了，那些被你们视为智者、被你们宣称是你们的宗教导师的人的混乱：他们中有些人宣称水是万物的第一原理；另一些人说是气；另一些人说是火；还有一些人说是前面提到的其他元素；所有人都用说服性的论证来建立自己的错误，并试图证明自己独特的教条是最有价值的。这些就是他们所说的话。那么，希腊人啊，对于那些渴望得救的人，怎能安全地幻想他们能从这些哲学家那里学到真正的宗教呢？这些哲学家既不能说服自己以致彼此之间没有派系争执，也不能不显得与彼此的观点明显对立。\par

% structure: p0024 source=h2 target=section reason=relative-heading-rank; base=h2

\section{柏拉图与亚里士多德的观点}

然而，那些不愿放弃古老而根深蒂固的错误之人，或许会坚持说，他们的宗教教义并非来自上述那些人，而是来自他们当中最负盛名、最完善的哲学家，即\Person{柏拉图}和\Person{亚里士多德}。因为他们说，这些人已学得完美而真正的宗教。但我想首先请问那些持此说者，他们认为这些人是向谁学得这种知识的？因为一个人若未从知道这些伟大而神圣之事的人那里学习，便不可能知道它们，也不可能正确地教导别人。其次，我认为我们应当审视这些智者的意见。因为我们将看到，他们每个人是否不公然彼此矛盾。但若发现他们甚至彼此不一致，我认为不难清楚地看出，他们也是无知的。因为\Person{柏拉图}带着从天而降、精确查明并看见了天上一切的神气，说至高天主存在于一种火性实体中。然而\Person{亚里士多德}在一本致\Place{马其顿}的\Person{亚历山大}的书中，简要解释了自己的哲学，清楚而明显地推翻了\Person{柏拉图}的意见，说天主不存在于火性实体中；而是发明了一种第五实体，即某种以太的、不变的身体，说天主存在于其中。他至少这样写道：\Quote{不像某些对神性有错误认识的人所说，天主存在于火性实体中。}然后，似乎不满足于对\Person{柏拉图}的这番亵渎，他为了证明自己关于以太之体的话，又引用了一个证人，即\Person{柏拉图}曾从他的共和国中驱逐的撒谎者，说他是模仿真理之像的第三级，因为\Person{柏拉图}这样称呼\Person{荷马}。他写道：\Quote{至少\Person{荷马}这么说过：\InnerQuote{宙斯获得了天空中广大的苍穹和云朵。}}他希望通过\Person{荷马}的见证使自己的意见显得更可信，却不知若用\Person{荷马}为证证明自己所言为真，他的许多主张就会被证明为假。因为\Place{米利都}的\Person{泰勒斯}，在他们当中是哲学的奠基人，会借机否定他关于第一原理的最初意见。因为\Person{亚里士多德}本人虽说过天主和物质是万物的第一原理，但所有智者中最年长的\Person{泰勒斯}却说水是万物的第一原理。他说万物源于水，又复归于水。他推测这一点，首先因为所有生物的种子（即它们的第一原理）是潮湿的；其次因为所有植物都在湿润中生长结果，失去湿润则枯萎。然后，似乎不满足于自己的推测，他引用\Person{荷马}作为最可信的见证，\Person{荷马}这样说：\par

\Quote{大洋，他是万物的起源。}\par

\Person{泰勒斯}岂非可以很公平地对他说：\Quote{\Person{亚里士多德}，当你想要推翻\Person{柏拉图}的意见时，为何相信\Person{荷马}，好像他说的是真理；而当你提出与我们相左的意见时，又认为\Person{荷马}不可信？}\par

% structure: p0028 source=h2 target=section reason=relative-heading-rank; base=h2

\section{第六章 \Person{柏拉图}与\Person{亚里士多德}的进一步分歧}

而你们这些非常奇妙的智者在其他方面也不一致，这一点从以下事实可以轻易看出：\Person{柏拉图}说万物的第一原理有三个：天主、物质和形式——天主是万物的创造者；物质是一切被造物最初产生的基质，为天主的制作提供可能性；形式则是每一个被造物的模型。而\Person{亚里士多德}根本不提形式作为第一原理，只说有两个：天主和物质。再者，\Person{柏拉图}说至高天主和理念存在于最高天界的固定领域；\Person{亚里士多德}却说，在至高天主之后，并非理念，而是某些能被心灵感知的神祇。因此，他们在天上之事上如此分歧。可见他们不仅不能理解我们地上的事物，而且在这些事上彼此矛盾，当他们论述天上之事时，也将显得不可信。而他们关于人类灵魂现今状态的学说也不一致，这一点从各自所说的话可以看出。因为\Person{柏拉图}说灵魂有三部分：理性、情感和欲望。但\Person{亚里士多德}说灵魂并非如此全面以至于包含可朽的部分，而只有理性。\Person{柏拉图}大声宣称\Quote{整个灵魂是不死的}。而\Person{亚里士多德}称它为\Quote{现实性}，认为它是有死的，而非不死的。前者说灵魂永远在运动；但\Person{亚里士多德}说灵魂是不动的，因为它本身必须先于一切运动。\par

% structure: p0030 source=h2 target=section reason=relative-heading-rank; base=h2

\section{第七章 \Person{柏拉图}学说的不一致之处}

但在这些事上，他们显然自相矛盾。若有人精确地批评他们的著作，他们甚至不愿与自己的见解保持一致。至少，\Person{柏拉图}有时说宇宙有三个本原——天主、物质与形式；但有时又说有四个，因为他加上宇宙灵魂。再者，当他已说过物质是永恒的，随后又说它是被造的；当他首先赋予形式作为本原的独特地位，并主张其自存性时，随后又说这同一事物属于理智所感知的事物。此外，他先宣称一切被造之物都是有死的，随后又说有些被造之物是不可毁灭、不朽的。那么，为何你们当中被视为有智慧的人不仅彼此不合，而且与自身也不合？显然，是因为他们不愿向认识者学习，却妄图凭藉自己过度的人类智慧获得对天上事物的精确认识，尽管他们连世间的事也无法理解。确实，你们的一些哲学家说人的灵魂在我们内，另一些说它在我们周围。就连这点他们也不愿彼此同意，反倒仿佛以各种方式在他们当中分配无知，认为应当就灵魂本身彼此争论和争辩。因为有些人说灵魂是火，有些说是空气，其他说是理智，其他说是运动，其他说是气息，还有一些说是从星宿流出的能力，其他说是能运动的数，其他说是生成之水。一种全然混乱而不和谐的意见在他们当中盛行，惟独在这一点上，对于那些能正确判断的人而言似乎值得称赞，即他们急切地要彼此证明错误与虚假。\par

% structure: p0032 source=h2 target=section reason=relative-heading-rank; base=h2

\section{第八章 基督徒导师的古老性、灵感与和谐}

既然因此不可能从你们的老师那里学到任何关于宗教的真理，因为他们彼此的歧见已向你们提供了充分的证据证明他们自己的无知，我认为合理的是回溯我们的祖先，他们不仅在时间上远远早于你们的老师，而且他们教导我们的并非出于任何私人的幻想，彼此也不相左，也不试图推翻对方的立场，而是毫无争执与争论，从天主那里接受了知识，并将之教导给我们。因为人既不能凭本性，也不能凭人的概念认识如此伟大而神圣的事物，而是靠从上往下赐予圣人们的恩赐，他们不需要修辞的技巧，也不需要以争论或吵闹的方式说话，而是要纯洁地呈现自己，以便接受圣神能力的感动，好使那神圣的拨子本身从天降下，以义人为乐器，如同竖琴或里拉琴，向我们启示神圣和天上事物的知识。因此，仿佛同一张嘴、同一舌，他们相继而和谐一致地教导了我们关于天主、世界的创造、人的形成、人类灵魂的不朽、此世之后的审判，以及一切我们需要知道的事物，并在不同的时期和地点赐予我们神圣的训诲。\par

% structure: p0034 source=h2 target=section reason=relative-heading-rank; base=h2

\section{第九章 由希腊作家证明\Person{梅瑟}的古老性}

那么，我要从我们的第一位先知和立法者梅瑟开始；首先依据在你们当中完全可信的权威，说明他所处的时代。因为我不仅打算从我们自己的神圣历史来证明这些事——由于你们祖先根深蒂固的谬误，你们至今不愿相信这些历史——也要从你们自己的历史中，以及那些与我们的敬拜毫无关系的著作中来证明，好使你们知道，在你们所有的导师中，无论是智者、诗人、历史学家、哲学家或立法者，按希腊历史所显示的，最为古老的是梅瑟，他是我们的第一位宗教导师。因为在奥巨格斯和伊纳科斯的时代——你们有些诗人认为他们是土生的——梅瑟就被提及为犹太民族的首领和统治者。因为波勒蒙在他的\WorkTitle{希腊史}第一卷中，以及波西多尼乌斯之子阿皮翁在他反对犹太人的书中，还有在他的历史第四卷中，都这样提到他：他说在伊纳科斯统治阿尔戈斯期间，犹太人从埃及王阿玛西斯统治下叛离，并由梅瑟领导他们。门德斯的托勒密在叙述埃及历史时，也同意这一切。而撰写雅典历史的赫拉尼库斯和菲洛科鲁斯（\WorkTitle{阿提卡史}的作者）、卡斯托尔、塔鲁斯、亚历山大·波利希斯托尔，以及非常了解犹太事务的斐罗和若瑟夫，都提到梅瑟是一位非常古老且备受尊崇的犹太君主。确实，若瑟夫甚至想通过他著作的标题来表明历史的古老和久远，他在历史著作的开头这样写道：\Quote{弗拉维乌斯·若瑟夫的\WorkTitle{犹太古史}}——用\Quote{古史}一词来表明历史的古老。你们最著名的历史学家狄奥多罗斯，他花费了整整三十年时间编写图书馆的摘要，并且如他自己所写，为了极大的准确性游历了亚洲和欧洲，从而亲眼目睹了许多事情，他编写了整整四十卷自己的历史。他在第一卷中说，他从埃及祭司那里得知梅瑟是一位古老的立法者，甚至是第一位，并用了这些话写到：\Quote{因为在埃及古老的生活方式——据传说由神祇和英雄所规范——之后，他们说是梅瑟首先说服民众使用成文法律并依此生活；他被记载为一位灵魂伟大且在社会事务上极具才能的人。}然后，他进一步叙述，并希望提到古老的立法者时，他首先提到了梅瑟。因为他这样说：\Quote{在犹太人当中，他们说梅瑟将他的法律归因于那被称为雅威的天主，这或许是因为他们认为这是一种奇妙且非常神圣的观念，承诺造福众多的人，或者因为他们认为，当民众仰望那些据说发明法律者的威严和能力时，会更加服从。他们说萨孙基斯是第二位埃及立法者，一位明智的人。第三位，他们说，是塞索恩科西斯王，他不仅在埃及做出了最辉煌的军事功绩，而且通过立法巩固了那个好战的民族。第四位立法者，他们说，是博科里斯王，一位明智且极为巧妙的人。之后，据说阿玛西斯王继位执政，他们讲述他规范了与地方统治者有关的一切事务，以及埃及政府的一般行政。他们说，薛西斯之父大流士是第六位为埃及人立法的人。}\par

% structure: p0036 source=h2 target=section reason=relative-heading-rank; base=h2

\section{第十章  梅瑟的培育与灵感}

希腊人啊，关于梅瑟的古老性，这些事是由那些不属于我们宗教的人写下的；他们说他们从埃及祭司那里得知了这一切。梅瑟不仅诞生在埃及人中，而且由于被法老女儿收养为子，被认为配得上分享埃及人的全部教育；由于同样的原因，他受到极大的重视，正如最明智的历史学家——我指的是斐罗和若瑟夫——所记载的，他们选择记录他的生活、事迹以及他的出身地位。因为他们在叙述犹太人的历史时说，梅瑟出自加色丁人的种族，并在埃及出生，当时他的祖先因饥荒从腓尼基迁移到那个国家；天主因他非凡的德行而拣选了他，并认为他配得上成为他本族的首领和立法者，那时天主认为希伯来人民应当从埃及返回自己的土地。天主首先将那天上降临在圣人们身上的神圣而预言的恩赐传达给梅瑟，也首先装备他成为我们宗教的导师，然后是他之后的其余先知，他们既获得了与他相同的恩赐，也在相同的主题上教导我们相同的教义。我们断言那些人是我们的导师，他们教导我们的不是任何来自自己人类构想的东西，而是来自从天上赐予他们的恩赐。\par

% structure: p0038 source=h2 target=section reason=relative-heading-rank; base=h2

\section{第十一章  外教神谕为梅瑟作证}

但既然你们认为服从这些（我们的）导师而放弃你们祖先的古老错误并无必要，那么你们主张在宗教问题上哪些你们自己的导师是值得信赖的呢？因为我已多次说过，那些没有从熟悉这些如此伟大、神圣之事的人那里学习过这些事的人，自己既不可能认识它们，也不可能正确地教导别人。既然如此，既然已充分证明你们哲学家的意见显然充满了一切无知和欺骗，现在也许像从前抛弃诗人一样彻底抛弃了哲学家，你们就会转向神谕的欺骗；因为我曾听一些人这样说过。因此，我认为在我们的论述中，此时告诉你们我以前在你们中间听到的关于他们神谕的回答是合适的。因为当有人询问你们的神谕——这是你们自己的故事——究竟何时曾有虔诚的人出现时，你们说神谕回答如下：\Quote{唯有加色丁人获得了智慧，以及希伯来人，他们崇拜天主本身，自生者君王。}\par

因此，既然你们认为真理可以从你们的神谕中得知，当你们阅读那些不属于我们宗教的人所写的历史和关于梅瑟生平的记载，并且知道梅瑟和其余先知都是加色丁人和希伯来人的后裔时，不要认为如果一个来自虔敬世系、并且生活与祖先的虔敬相称的人被天主拣选，蒙受如此伟大的恩赐，并被立为所有先知中的首位，是有什么不可思议的事。\par

% structure: p0041 source=h2 target=section reason=relative-heading-rank; base=h2

\section{第十二章 梅瑟的古老性之证明}

我认为也有必要考察一下你们哲学家的年代，以使你们看到他们为你们产生的时代是非常晚近且短暂的。因为这样你们就能容易地认识到梅瑟的古老性。但为了避免通过全面考察各个时期和使用更多证据而显得冗长，我认为从以下几点就足以证明。因为苏格拉底是柏拉图的老师，柏拉图是亚里士多德的老师。这些人活跃于马其顿的斐理伯和亚历山大时期，在这个时期，雅典的演说家们也活跃着，正如德摩斯梯尼的\WorkTitle{斥斐力辞}清楚地告诉我们的。那些叙述亚历山大事迹的人也充分证明，在他的统治期间，亚里士多德曾与他交往。因此，从各种证据来看，很容易看出梅瑟的历史远早于所有世俗历史。此外，你们也应当认识到这一点：在奥林匹克纪元之前，希腊人没有精确记载任何事物，也没有任何古代著作记载希腊人或蛮族的任何行动。但在那个时期之前，只有先知梅瑟的历史存在，他是因天主默感而用希伯来文字写成的。因为当时希腊文字尚未使用，正如语言教师们自己所证明的，他们告诉我们，卡德摩斯首先从腓尼基带来了字母，并传给了希腊人。你们的首位哲学家柏拉图也证实它们是一项新近的发现。因为他在\WorkTitle{蒂迈欧篇}中写道：梭伦，智者们中最智慧的一位，从埃及回来后对克里底亚说，他从一位非常年迈的埃及祭司那里听说了这事，那位祭司对他说：\Quote{啊，梭伦，梭伦，你们希腊人永远是孩童，没有年迈的希腊人。}然后他又说：\Quote{你们在灵魂上都是幼童，因为你们不拥有通过遥远传统传承的任何古老意见，也没有任何因时间久远而灰白的教训体系；但所有这些事你们都不知道，因为许多世代以来，这些古老时代的后代尚未使用文字就默默死去了。}因此，你们应当明白，事实上每一部历史都是用这些新近发现的希腊文字写成的；如果有人想提及古老的诗人、立法者、历史学家、哲学家或演说家，他会发现他们是用希腊文字写成自己的作品的。\par

% structure: p0043 source=h2 target=section reason=relative-heading-rank; base=h2

\section{第十三章 七十贤士译本史}

但若有人声称\WorkTitle{梅瑟}和其余先知的作品也是用希腊文字写成的，就让他去读世俗的历史吧，并且要知道：埃及王\Person{托勒密}在\Place{亚历山大}建成了图书馆，从各处收集书籍填满之后，得知有一些用希伯来文写成的极古老的历史著作被小心保存着；他渴望知道其内容，便从\Place{耶路撒冷}请来了七十位精通希腊语和希伯来语的贤士，任命他们翻译这些书；为了使他们在不受任何干扰的情况下更快地完成翻译，他命令在城外（\Place{法罗斯}灯塔所在地）七斯塔狄昂处建造与译员数量相同的小屋，以便每人独自完成自己的翻译；并命令负责此事的官员向他们提供一切服务，但阻止他们彼此交流，以便通过他们的一致来辨别翻译的准确性。当他确定这七十人不仅给出了相同的意思，而且使用了相同的词语，彼此之间连一个字的差异都没有，而是写了相同的事物、关于相同的事物时，他大为惊异，相信这译本是凭藉天主的能力写成的，并看出这些人是值得一切尊荣的，如同天主所爱的人；于是他赐予许多礼物，命他们返回本国。他自然对这些书惊叹，断定它们是神圣的，便将它们奉藏在图书馆中。这些事情，希腊人啊，并非寓言，我们也不是在编造虚构的故事；我们本人在\Place{亚历山大}时，亲眼见过\Place{法罗斯}那里那些小屋的遗迹至今仍保存着，并从当地居民——他们将这些作为本国传统的一部分传承下来——听说了这些事，现在我们将这些告诉你们，你们也能从其他人那里得知，尤其是从那些明智而可敬的、曾记述这些事的人，如\Person{斐罗}、\Person{若瑟夫}和许多其他人。但若有人惯于反驳，说这些书不属于我们，而属于犹太人，并断言我们自称从他们那里学到我们的宗教是徒然的，那就让他从这些书中所写的内容本身知道：它们的教义不是指向他们，而是指向我们。与我们的宗教相关的书至今保存在犹太人中间，这乃是对我们而言的天主眷顾之举；因为，若我们从教会内拿出这些书，会给那些想要诽谤我们的人以借口，指责我们欺诈，所以我们要求从犹太人的会堂中拿出它们，以便通过仍保存在他们当中的这些书本身，清楚明白地显明：由圣人为教导而写下的法律是属于我们的。\par

% structure: p0045 source=h2 target=section reason=relative-heading-rank; base=h2

\section{第十四章 对希腊人的警告呼吁}

因此，希腊人啊，你们必须默观将要发生的事，并思考那由所有人——不仅由虔诚的人，也由不敬神的人——所预言的审判，以免未经考察就陷入你们父辈的错误，也不要以为他们自己既然犯了错并将之传给你们，他们所教导你们的便是真理；而要看到如此可怕错误的危险，仔细探究和考察那些连你们的老师自己都曾说过的事。因为即使是不情愿地，出于天主对人类的眷顾，他们也被迫为你们说了许多事，尤其是那些曾在\Place{埃及}并得益于\Person{梅瑟}及其祖先的虔诚的人。因为我认为，你们中有些人，即使漫不经心地阅读\Person{狄奥多罗斯}和那些记述这些事之人的历史，也不能不看到：无论是\Person{俄耳甫斯}、\Person{荷马}、撰写\Place{雅典}法律的\Person{梭伦}，还是\Person{毕达哥拉斯}、\Person{柏拉图}及其他一些人，当他们到过\Place{埃及}并利用过\Person{梅瑟}的历史之后，后来发表的关于诸神的学说，与他们先前错误传播的完全不同。\par

% structure: p0047 source=h2 target=section reason=relative-heading-rank; base=h2

\section{第十五章 \Person{俄耳甫斯}对一神论的见证}

无论如何，我们必须提醒你们：那位可以说是你们多神论的首位导师的\Person{俄耳甫斯}，后来对他的儿子\Person{穆赛俄斯}和其他合法的听众所说的话，关于那唯一的一位天主。他这样说道：\par

\Quote{我对那些合法倾听的人说话：\newline{}其余人等，亵渎者啊，现在关门吧，\newline{}哦，\Person{穆赛俄斯}！请听我说，\newline{}你是光耀之月的后裔：\newline{}我现在所说的话实属真实；\newline{}若你曾见过我以往的想法，\newline{}不要让它们夺去你蒙福的生命，\newline{}而要将你内心深处的意念转向\newline{}那光明与智慧居住之处。\newline{}以神圣的话引导你的步伐，\newline{}行在笔直确定的道路上，\newline{}仰望那唯一的宇宙之王——\newline{}唯一，自生，且独一无二，\newline{}万物及我们自己都由祂而生。\newline{}万物都为祂锐利的目光所洞悉，\newline{}而祂自身却依然不可见。\newline{}在祂所有作品中呈现，仍隐而不现，\newline{}祂从善中赐给凡人恶，\newline{}降下冷酷的战争和悲伤的哭泣；\newline{}除那伟大的王之外别无他者。\newline{}云雾永远环绕祂的宝座，\newline{}凡人眼目在纯粹肉眼中\newline{}是软弱无力的，无法看见统领万有的\Person{宙斯}。\newline{}祂安坐在青铜天上\newline{}的金色宝座上；祂的脚下\newline{}踩着大地，右手伸展\newline{}到海洋的尽头，四周\newline{}群山颤抖，溪流战栗，\newline{}青灰大海的深处也一同震动。}\par

在另一处，他又说：\par

\Quote{只有一个\Person{宙斯}，一个太阳，一个地狱，\newline{}一个\Person{巴克科斯}；万物中只有一位天主；\newline{}且我不要说这一切是各不相同的。}\par

当他发誓时，他说：\par

\Quote{现在我以最高的天向你起誓，\newline{}那是伟大天主、唯一智慧者的作品；\newline{}我又以圣父的声音向你起誓，\newline{}祂在凭自己的旨意建立\newline{}整个世界时首先发出了那声音。}\par

他所说的\Quote{我指着圣父最先发出的声音命令你}是什么意思？他在这里称天主的圣言为\Quote{声音}，藉此声音天地及万有被造，正如圣人们的神圣预言教导我们的；他自己在埃及也曾留意这些预言，并理解万有是由天主的圣言所造；因此，在他说\Quote{我指着圣父最先发出的声音命令你}之后，他又加上\Quote{当祂以祂的计谋建立了整个世界}一句。在这里他为了诗韵的缘故称圣言为\Quote{声音}。这一点是显而易见的，因为在稍后，当韵律允许时，他便称之为\Quote{圣言}。因为他写道：\par

\Quote{以神圣的圣言指引你的步履。}\par

% structure: p0056 source=h2 target=section reason=relative-heading-rank; base=h2

\section{第十六章　西彼尔的见证}

我们也必须提及那位古老而极为遥远的西彼尔，\Person{柏拉图}、\Person{阿里斯托芬}以及其他人均提到她是一位女先知，她以神谕诗篇教导你们关于独一天主的道理。她这样说：——\par

\Quote{有一位独一非受生的天主，\newline{}全能、不可见、至崇高，\newline{}洞察万物，却无血肉之躯得见祂。}\par

接着在别处这样说：——\par

\Quote{但我们已偏离永生者的道路，\newline{}以迟钝无知的心灵崇拜\newline{}偶像，我们双手的制造，\newline{}以及死人的形象和雕像。}\par

又在某处说：——\par

\Quote{有福的是地上那些人，\newline{}他们必先爱伟大的天主，\newline{}饮食之时都祝福祂；\newline{}唯独信赖这敬虔。\newline{}他们必弃绝所见的一切庙宇，\newline{}所有祭台和哑巴石头的虚妄形象，\newline{}毫无价值且沾染兽血，\newline{}以及四足族类的祭献，\newline{}瞻仰独一天主的伟大荣耀。}\par

以上是西彼尔的话语。\par

% structure: p0064 source=h2 target=section reason=relative-heading-rank; base=h2

\section{第十七章　荷马的见证}

诗人\Person{荷马}运用诗的许可，与\Person{俄耳甫斯}关于诸神多元的原本观点竞争，确实以神话风格提及多位神祇，以免他显得与俄耳甫斯的诗篇风格不同，而他显然立志与之竞争，以致在他的诗篇第一行就表明了他与俄耳甫斯的关系。因为俄耳甫斯在其诗篇开头曾说：\Quote{女神啊，请歌唱带来佳果的\Person{德墨忒耳}的愤怒。}荷马则这样开始：\Quote{女神啊，请歌唱\Person{珀琉斯}之子\Person{阿喀琉斯}的愤怒。}在我看来，他甚至宁愿在第一行破坏诗韵，也不愿显得没有首先记住神祇的名字。但不久之后，他也清楚而明确地提出了自己关于独一天主的见解，在某处借\Person{福尼克斯}之口对阿喀琉斯说：\Quote{即使天主亲自应许褪去我的老年，还我青春的活力，}他在这里用指示代词指出了真实和真正的天主。又在某处他让\Person{尤利西斯}对希腊军队这样说：\Quote{多人统治并非好事；让一人为统治者。}至于多人统治并非好事，反而是恶，他打算用事实来证明，叙述了因众多统治者而发生的战争、争斗、派系以及相互的阴谋。因为君主政体\Emph{是}免于纷争的。诗人荷马这样说。\par

% structure: p0066 source=h2 target=section reason=relative-heading-rank; base=h2

\section{第十八章　索福克勒斯的见证}

若有必要，我们甚至从戏剧家那里添加关于独一天主的见证，请听\Person{索福克勒斯}这样说：——\par

\Quote{有一位天主，真理中只有一位，\newline{}祂造了天和下方广阔的大地，\newline{}闪亮的海浪与风；\newline{}但我们许多有死者心中犯错，\newline{}为慰藉我们的苦难\newline{}立起石木制成的神像，\newline{}或用铜或象牙雕成的形象，\newline{}并为这些我们的手工品\newline{}献上辉煌的祭献和仪式，\newline{}以为这样做便是虔敬之举。}\par

以上是索福克勒斯的话。\par

% structure: p0070 source=h2 target=section reason=relative-heading-rank; base=h2

\section{第十九章　毕达哥拉斯的见证}

\Person{莫涅撒库斯}之子\Person{毕达哥拉斯}，用象征方式神秘地阐述他自己哲学的教义，如那些为他作传的人所示，他本人看似对天主的合一怀有见解，与其在\Place{埃及}的异邦居留不相称。因为当他说统一是万物的第一原理，且是一切美善的原因时，他通过寓言教导天主是独一且唯一的。这一点从他所说统一与一彼此大相径庭也可明显看出。因为他说统一属于心智所感知的事物之类，而一属于数目。你若想看到关于独一天主的毕达哥拉斯观点更清晰的证明，请听他的见解，因为他这样说：\Quote{天主是独一的；祂并非如有些人所设想的，存在于世界之外，而是在世界之内，祂完全临在于整个圆环中，并洞悉一切世代；祂是万有的调节成分，是祂自身能力和工作的管理者，是万物的第一原理，是天上的光，是万有之父，是宇宙的理智和赋灵的灵魂，是一切轨道的运动。}以上是毕达哥拉斯的话。\par

% structure: p0072 source=h2 target=section reason=relative-heading-rank; base=h2

\section{第二十章　柏拉图的见证}

但柏拉图虽然很可能接受了梅瑟及其他先知关于独一真天主的教导——这是他在埃及学到的——然而由于害怕苏格拉底的遭遇，唯恐他自己也会招来某个阿尼托或梅勒托在雅典人面前控告他，说：\Quote{柏拉图在作恶，多管闲事，不承认国家所承认的神明；}出于对毒芹汁的恐惧，他构思了一篇关于诸神的精巧而模糊的论述，通过他的论文为那些希望有神的人提供神，而为那些持不同态度的人则没有，这从他自己的陈述中可以轻易看出。因为他既已规定一切受造的都是有死的，此后又说诸神是被造的。那么，如果他主张天主和物质是万物的本原，显然就不可避免地必须说诸神是由物质造成的；但如果是出自物质——他说邪恶也源于物质——他让思想正确的人去思考那些由物质产生出来的神应当被看作什么样的存在者。因为正是为此，他说物质是永恒的，以免他显得说天主是邪恶的创造者。至于那些由天主所造的诸神，无疑他说过：\Quote{诸神中的诸神，我是它们的创造者。}关于真正存在的天主，他显然持有正确的意见。因为他在埃及听说，当天主即将派遣梅瑟到希伯来人那里时，对梅瑟说：\Quote{我是自有者，}他就明白天主没有向梅瑟提及祂自己的专有名称。\par

% structure: p0074 source=h2 target=section reason=relative-heading-rank; base=h2

\section{第二十一章 天主的无名性}

因为任何人都不能给天主起一个专有名称，因为名称是用来标示和区分其主体，而这些主体是众多且多样的；但在天主之前并没有任何人存在能给祂起名，祂自己也认为不当为自己命名，因为祂是唯一无二的，正如祂自己藉祂的先知所见证的，当祂说：\Quote{我天主是首先的，}并在此之后：\Quote{除我以外没有别的神。} \BibleRef{依 44:6} 因此，正如我之前所说，天主在派遣梅瑟到希伯来人那里时，并没有提及任何名字，而是用一个分词神秘地教导他们祂是唯一的天主。祂说：\Quote{因为，} \Quote{我是自有者；}显然将\Quote{自有者，}与那些不存在者相对比，以便那些迄今受骗的人能够看到他们依附的不是存在者，而是那些没有存在的东西。因此，既然天主知道最初的人还记得他们祖先的古老幻象，那厌恨人类的魔鬼藉此设计欺骗他们，他对他们说：\Quote{如果你们听从我违犯天主的诫命，你们将如同诸神，}他把那些不存在的东西称为诸神，为使人以为存在其他神，从而相信他们自己也能成为神。为此，他对梅瑟说：\Quote{我是自有者，}以便通过分词\Quote{存有}教导祂存在的天主与那些不存在者之间的区别。因此，人受了欺骗魔鬼的愚弄，胆敢违抗天主的命令，被逐出乐园，他们还记得诸神的名字，但不再受天主教导说没有别的神。因为对于那些没有遵守第一条容易遵守的诫命的人，不再适合教导他们，而应遭受公义的惩罚。因此，他们被逐出乐园，以为只是因违命被驱逐，不知道也是因为他们相信了并不存在的诸神的存在，他们甚至把他们自己后出生的人也称作神。因此，关于诸神的第一个错误幻想起源于说谎之父。因此，天主知道关于多神的错误观念像某种疾病一样重压着人的灵魂，希望除去并根除它，就首先向梅瑟显现，对他说：\Quote{我是自有者。}因为我认为，那将要作希伯来民族统治者和领袖的人必须先认识生活的天主。所以，天主首先向他显现，正如天主向人显现可能的那样，祂对他说：\Quote{我是自有者；}然后，即将派遣他到希伯来人那里时，又命令他对他们说：\Quote{那自有者派遣我到你们这里来。}\par

% structure: p0076 source=h2 target=section reason=relative-heading-rank; base=h2

\section{第二十二章 刻意模糊的柏拉图}

\Person{柏拉图}在\Place{埃及}学到了这一点，并且对那关于独一真天主的言论大为倾倒，但鉴于\Person{梅瑟}的训诲只教导一位天主，他认为提及\Person{梅瑟}的名字并不安全，因为他惧怕\Place{阿勒约帕哥}；然而，在他精心撰写的论著\WorkTitle{蒂迈欧篇}中，他所写的与\Person{梅瑟}关于天主的论述完全吻合，尽管他这样做，并非像是从\Person{梅瑟}那里学来的，而是像是在表达他自己的见解。因为他说：\Quote{\Emph{在我看来}，那么我们必须首先界定那永恒存在、没有生成的是什么，以及那总是被生成、却从未真正存在的是什么。}各位希腊人啊，这对那些能领会此事的人来说，岂不像是同一个东西，只是冠词的差别吗？因为\Person{梅瑟}说：\Quote{\Emph{祂}是那自有者}，而\Person{柏拉图}说：\Quote{那是者}。但这两个表述似乎都指向永远存在的那位天主。因为祂是唯一永恒存在、没有生成的。那么，与那永恒存在者相对的另一物是什么，他所说的\Quote{那总是被生成、却从未真正存在的是什么}，我们必须仔细思考。因为我们会发现他清楚明确地说，那非受生的是永恒的，但那受生和被造的则被生成而毁灭——正如他论及同一类所说：\Quote{众神之神，我是它们的制造者}——因为他这样说道：\Quote{\Emph{在我看来}，那么我们必须首先界定那总是存在、没有出生的东西是什么，以及那总是被生成、却从未真正存在的东西是什么。前者，的确，由反思与理性所把握，总是以同样的方式存在；而后者，则是由感官知觉形成的意见所推测，没有理性，因为它从未真正存在，而是在生成和毁灭中。}这些话向那些能正确理解的人宣告了那些被引入存在的众神的死亡与毁灭。而且我认为有必要注意这一点：\Person{柏拉图}从未称他为创造者，而是称他为众神的塑造者，尽管在\Person{柏拉图}看来，这两者之间有相当大的区别。因为创造者凭自身的能力和权能创造受造物，无需其他东西；但塑造者当他从物质那里获得了他作品的能力时，才塑造他的作品。\par

% structure: p0078 source=h2 target=section reason=relative-heading-rank; base=h2

\section{第二十三章　\Person{柏拉图}的自相矛盾}

但是，也许有些人不愿意放弃多神论的教义，他们会说，制造者对那些被塑造的众神说：\Quote{既然你们已被产生，你们并非不朽，也绝非不坏；但你们不会毁灭，也不会屈服于死亡的命运，因为你们获得了我的意志，这是一个更大更强有力的纽带。}在这里，\Person{柏拉图}由于害怕多神论的追随者，让他的\Quote{制造者}说出与自己矛盾的话。因为他先前曾说他断言一切被产生的都是可朽的，现在却让他说出完全相反的话；而且他没有看到，这样他绝对不可能逃脱说谎的指控。因为他要么最初说的就是假话，当他说一切被产生的都是可朽的；要么现在，当他提出与先前所说的完全相反时。因为如果根据他最初的定义，每一个被造物都必然是可朽的，那么他如何能前后一致地使那绝对不可能的事成为可能呢？所以当他提出那些曾经因由物质构成而是可朽的，能因他的介入而再次成为不朽和持久的时，\Person{柏拉图}似乎赋予他的\Quote{制造者}一个空洞且不可能的特权。因为很自然，物质的力量（按\Person{柏拉图}的观点，它是非受造的，与制造者同时代且同样永恒）会抵抗他的意志。因为那没有创造的人，对于那非受造的，没有权能，以至于它（物质）既然是自由的，就不可能被任何外在的必然性所控制。因此，\Person{柏拉图}自己考虑到这一点，这样写道：\Quote{必须肯定，天主不能受强暴。}\par

% structure: p0080 source=h2 target=section reason=relative-heading-rank; base=h2

\section{第二十四章　\Person{柏拉图}与\Person{荷马}的一致}

那么，\Person{柏拉图}如何将\Person{荷马}从他的理想国中驱逐出去呢？因为在向\Person{阿基里斯}的使团中，他让\Person{腓尼克斯}对\Person{阿基里斯}说：\Quote{即使是众神自己也并非不可改变}，虽然\Person{荷马}这样说不是指那位众神之王和\Person{柏拉图}式的制造者，而是指那些被希腊人尊为神的一些群众，正如可以从\Person{柏拉图}的话\Quote{众神之神}中推断的那样。因为\Person{荷马}通过那根金链，将一切权能和能力归于那至高的独一真神。至于其余的神，他说他们离他的神性如此遥远，以至于他认为甚至提到他们时也可以与人并列。至少他让\Person{尤利西斯}谈到\Person{赫克托耳}对\Person{阿基里斯}时说：\Quote{他狂暴地肆虐，信赖\Person{宙斯}，既不重视人也不重视神。}在这一段中，\Person{荷马}在我看来毫无疑问地像\Person{柏拉图}一样在\Place{埃及}学到了关于独一真神的道理，并且明白公开地宣称：那信赖真正存在之天主的人，不会顾念那些不存在的东西。因为诗人在另一段中，用了另一个却同等的词语，即一个代词，使用了与\Person{柏拉图}相同的分词来指代真正存在之天主，关于这一点\Person{柏拉图}曾说过：\Quote{那总是存在、没有出生的东西是什么。}因为\Person{腓尼克斯}的这句话似乎不是无意被使用的：\Quote{即使天主\Emph{亲自}应许我，磨去我的老迈，使我重返青春。}代词\Quote{亲自}指代那真正存在之天主。因为这样，那曾给予你们关于迦勒底人和希伯来人的神谕也表明。当有人问起，曾经有哪些人虔诚地生活过，你们说回答是：\par

\Quote{惟独迦勒底人和希伯来人寻得智慧，\newline{}敬拜天主亲自，那非受生的君王。}\par

% structure: p0083 source=h2 target=section reason=relative-heading-rank; base=h2

\section{第二十五章　\Person{柏拉图}对天主永恒的认识}

那么，柏拉图为何责备荷马说众神并非坚定不移，尽管从所用的措辞来看，荷马这样说是有益的目的呢？因为凡期望藉祈祷和祭献获得仁慈的人，其特性就是停止并悔改自己的罪。因为那些认为神明不可改变的人，绝不会被推动去放弃他们的罪，因为他们以为从悔改中得不到任何益处。那么，哲学家柏拉图如何能责备诗人荷马说\Quote{连众神本身也并非坚定不移，}而他自己却把创造众神的神明描述得如此轻易转变，有时宣称众神是有死的，有时又宣称它们是不死的？而且不仅关于它们，而且关于质料（据他说，被造的神必须从质料中产生），他有时说它是非受造的，有时又说它是受造的；然而他却没有看到，当他说创造众神的神明如此轻易转变时，他自己就被证实陷入了他曾责备荷马的同样错误，尽管荷马关于创造众神的神明所说的正相反。因为他说，荷马是这样说到自己的：——\par

\Quote{因为我的诺言永远不会欺骗、落空，\newline{}或撤回，若以点头确认。}\par

但是柏拉图似乎不情愿地加入了这些关于诸神的奇怪讨论，因为他害怕那些依附多神论的人。而他想到从梅瑟和先知那里学到有关唯一天主的任何事，他喜欢以神秘的方式传授，以便那些渴望成为天主崇拜者的人能对他的意见有所察觉。因为他被天主对梅瑟所说的那句话\Quote{我是那真实存在的，}所吸引，并带着大量的思考接受了那个简短的分词表达，他明白天主想向梅瑟表示祂的永恒，因此说\Quote{我是那真实存在的；}因为\InnerQuote{存在}这个词不仅表达一个时间，而是三个——过去、现在和未来。因为当柏拉图说\Quote{且从未真实存在，}他用动词\Emph{是}表示不确定的时间。因为\Emph{从未}一词并非如某些人所想，指的是过去，而是将来。这一点甚至被世俗作家准确地理解了。因此，当柏拉图想要为未入门者解释那分词形式所神秘表达的关于天主永恒性时，他使用了以下语言：\Quote{天主，正如古老的传统所说，包含万物的开始、终结和中间。}在这句话中，他明白而明显地将梅瑟的法律称为\Quote{古老的传统，}因畏惧毒芹杯而不敢提梅瑟的名字；因为他明白那人的教导为希腊人所憎恨；而他通过传统的古老性足够清楚地指明了梅瑟。我们在前几章中从狄奥多洛斯和其他历史学家那里已经充分证明，梅瑟的法律不仅是古老的，而且甚至是第一部。因为狄奥多洛斯说他是所有立法者中的第一个；属于希腊人的文字，他们用来撰写历史的，还没有被发现。\par

% structure: p0087 source=h2 target=section reason=relative-heading-rank; base=h2

\section{第26章　柏拉图对先知负有亏欠}

任何人都不应奇怪柏拉图相信梅瑟关于天主永恒性的说法。因为你会发现在他神秘地将关于现实界的真正知识归诸先知，仅次于那真实存在的天主之后。因为，在\WorkTitle{蒂迈欧篇}中讨论某些第一原理时，他写道：\Quote{我们确立这作为火和其他物体的第一原理，按概然性和必然性进行。但这些的第一原则，唯有天主知道，以及凡蒙祂所爱的人。}那么，他认为哪些人是蒙天主所爱的呢？难道不是梅瑟和其余的先知吗？因为他读了他们的预言，并从他们学到了审判的道理，于是在\WorkTitle{共和国}的第一卷中这样宣告：\Quote{当一个人开始想到自己很快就要死时，恐惧就侵入他，担忧那些从未进入他头脑的事物。那些关于冥府中所发生的事的故事，告诉我们今世不义的人必须在彼世受惩罚，尽管从前被嘲笑，现在却以担忧它们可能是真的而折磨他的灵魂。他，或因年老的衰弱，或因其更接近另一个世界的事物，便更专注地看待它们。于是他充满忧虑和恐惧，开始反躬自问，考虑是否曾对任何人造成伤害。那发现自己一生中有许多不义的人，像儿童一样从睡梦中惊醒，生活在恐惧中，前途凄惨。但自觉无过的人，甜美的希望常伴左右，是老年良好的看护，如品达所说。因为，苏格拉底，他优美地表达了：\InnerQuote{凡度圣洁与正义生活的人，甜美希望，衰老的看护，伴随着他，愉悦他的心灵，因为她有力地支配着凡人变动的心。}}柏拉图在\WorkTitle{共和国}第一卷中就是这样写的。\par

% structure: p0089 source=h2 target=section reason=relative-heading-rank; base=h2

\section{第27章　柏拉图对审判的认识}

在第十卷中，他清晰而明确地写下了他从先知们那里学到的关于审判的教导，但不是以从他们那里学来的方式，而是由于他对希腊人的恐惧，仿佛他是从一个在战斗中阵亡的人那里听来的——他认为编造这个故事是合适的——那个人在十二天后将要下葬时，躺在火葬堆上，复活了，并描述了另一个世界。以下是他原话：\Quote{因为他说，他曾亲眼目睹有人问另一个人，伟大的阿尔迪乌斯在哪里。这个阿尔迪乌斯曾是旁非利亚某城的王子，他杀死了年迈的父亲和长兄，并犯下了许多其他不圣洁的罪行，如传闻所说。他接着说，被问的人回答：\InnerQuote{他既不來，也永远不会来这里。因为我们看到，在其他可怕的景象中，还有这一件。当我们靠近（深渊的）口，正要返回上界，并经历了其他一切苦难时，我们突然看见了他，以及另外一些人，其中大多是暴君。但也有几个犯下大罪的私人罪人。这些人在他们以为自己要上升时，那口却不允许，而是当那些无可救药的恶人试图上升时，除非他们受到完全的惩罚，否则口会发出吼声。然后有一些凶猛的人，看起来像火焰一样，站在旁边，听到喧闹声，就抓住一些人并带走他们；但阿尔迪乌斯和其余的人，手脚被捆绑，头被击打，皮被剥去，他们被拖到外面的路上，用荆棘撕裂，并向在场的人表明他们遭受这些苦难的原因，以及他们正在被带去投入塔尔塔罗斯。}因此，他说，在他们各种恐惧中，最大的恐惧是当他们上升时口会发出吼声，因为如果口沉默，每一个人都会最乐意上升；而惩罚和折磨就是如此，另一方面，奖赏则与此相反。} 在我看来，柏拉图不仅从先知们那里学到了审判的教义，还学到了希腊人所拒绝相信的复活教义。因为他所说的灵魂与身体一同受审判，无非是证明他相信复活的教义。因为阿尔迪乌斯和其余的人怎么能在哈得斯遭受这样的惩罚，如果他们把身体连同头、手、脚和皮肤都留在地上？当然，他们绝不会说灵魂有头、手、脚和皮肤。但柏拉图在埃及遇见了先知们的见证，并接受了其中关于身体复活的教导，因此他教导说灵魂与身体一同受审判。\par

% structure: p0091 source=h2 target=section reason=relative-heading-rank; base=h2

\section{第二十八章 荷马对圣经作者的借鉴}

不仅柏拉图，而且荷马也在埃及得到了类似的启示，说提堤俄斯受到同样的惩罚。因为尤利西斯在向阿尔基努斯讲述他通过死者的阴魂所进行的占卜时，说了这样的话：\par

\Quote{那里提堤俄斯，巨大而修长，被锁链捆绑，\newline{} 覆盖了九英亩的地狱地面；\newline{} 两只贪婪的秃鹫，因食物而狂怒，\newline{} 在那恶魔上方尖叫，并在他的血中狂欢，\newline{} 不停地啄食他胸中的肝脏，\newline{} 不朽的肝脏生长，给予不朽的盛宴。}\par

因为显然，有肝脏的不是灵魂，而是身体。同样，他也描述了西绪福斯和坦塔罗斯以身体承受惩罚。而且荷马曾到过埃及，并将他在那里学到的许多内容引入自己的诗歌中，这一点最受尊敬的史学家狄奥多罗斯就明白地教导我们。因为他说，当他在埃及时，他得知海伦从特翁的妻子波吕达姆娜那里获得了一种药物，\Quote{消除一切忧伤和忧愁，并使人忘却所有苦难}，并将它带到了斯巴达。荷马说，海伦使用这种药物，平息了墨涅拉俄斯因忒勒马科斯在场而引起的悲伤。他还从他在埃及所见而称维纳斯为\Quote{黄金的}。因为他看见在埃及被称为\Quote{黄金维纳斯殿}的庙宇，以及被称为\Quote{黄金维纳斯平原}的平原。我现在为何提起这些？是为了表明诗人将他从先知神圣著作中获得的许多内容转移到了自己的诗中。首先，他转移了摩西关于世界创造开端的叙述。因为摩西这样写道：\Quote{起初天主创造了天地}，\BibleRef{创 1:1}，然后是太阳、月亮和众星。因为他在埃及学到了这些，并且对摩西在\BibleBook{}{创世记}中所写的关于世界创造的内容非常着迷，于是他虚构说，武尔坎在阿喀琉斯的盾牌上制作了类似世界创造的一种描述。因为他这样写道：\par

\Quote{他在那里描绘了大地、天空、海洋，\newline{} 不息的太阳和满盈的月亮；\newline{} 在那里，还有所有星辰，\newline{} 如同灿烂的额带，环绕着天空。}\par

他还构想出阿尔基努斯的花园，使之保存了乐园的相似，并通过这种相似将其描绘为永远开花、充满各种果实。因为他这样写道：\par

\Quote{高大的树木欣欣向荣，承认肥沃的土壤；\newline{} 泛红的苹果成熟变成金色。\newline{} 在这里，蓝色的无花果洋溢着甜美的汁液，\newline{} 更深红的石榴闪耀着红光；\newline{} 树枝在这里因沉重的梨子而低垂，\newline{} 肥沃的橄榄一年四季繁茂。\newline{} 西风的芳香气息\newline{} 永远吹拂着果实，永不凋零；\newline{} 每个掉落的梨子都有下一个梨子接续，\newline{} 苹果之上又生苹果，无花果之上又生无花果。\newline{} 同样的温和季节使花朵绽放，\newline{} 使花蕾结实，使果实生长。\newline{} 这里排列整齐的葡萄藤以相等的行列出现，\newline{} 伴随着一年中所有联合的劳动。\newline{} 有的人跑去卸下丰饶的树枝上的果实，\newline{} 有的人在阳光下晒干变黑的葡萄串，\newline{} 其他人则加入踩踏液体的收获。\newline{} 呻吟的压榨机喷涌出葡萄酒的洪流。\newline{} 这里可见葡萄藤在早花中，\newline{} 这里葡萄在阳光照耀下变色，\newline{} 而在那里，则在秋日最深的紫色中染色。}\par

这些话难道不是对第一位先知\Person{梅瑟}关于乐园所说的话的明显而清晰的模仿吗？如果有人想了解当时的人妄想以此登天的那座塔的建造，他会在诗人归于\Person{奥托斯}和\Person{厄菲阿尔特斯}的描写中找到足够精确的寓意模仿。因为诗人这样写道：\par

\Quote{他们自恃力大无比，超乎凡人，\newline{}竟敢挑战众神，觊觎天庭。\newline{}他们将\Place{奥萨山}摇摇欲坠地堆在\Place{奥林波斯山}上，\newline{}又在\Place{奥萨山}上堆起\Place{佩利昂山}，林木茂密。}\par

关于那从天上被逐出的人类之敌也是如此，圣经称他为魔鬼，这名称源于他最初对人类的恶行；如果任何人仔细思考此事，就会发现诗人虽从未提及\Quote{魔鬼}之名，却以他最邪恶的行为给他起了名字。因为诗人称他为\Person{阿忒}，说他被他们的神从天上扔下，仿佛他清楚记得先知\Person{依撒意亚}关于他所说的话。诗人在他的诗中这样写道：\par

\Quote{他抓住女神\Person{阿忒}油亮的发辫，\newline{}盛怒之下起誓，\newline{}决不再让这普世的祸害\newline{}返回星空和\Place{奥林波斯山}巅。\newline{}然后他旋转她，\newline{}将她抛向大地。\newline{}她混入人的一切作为，\newline{}给\Person{宙斯}带来许多痛苦。}\par

% structure: p0102 source=h2 target=section reason=relative-heading-rank; base=h2

\section{第二十九章 柏拉图理型学说的起源}

而\Person{柏拉图}说\OriginalTerm{理型}{form}是次于天主和物质的第三个本原时，显然也是从\Person{梅瑟}那里得到这一暗示；他确实从\Person{梅瑟}的话中学到了理型的名称，但未同时受到入门者的教导，即没有神秘洞见就不可能清楚理解\Person{梅瑟}的著作。因为\Person{梅瑟}记载天主关于会幕对他说的话如下：\BibleQuote{你要按照我在山上向你展示的，为我建造会幕的样式。}\EditorialNote{出谷纪 25}又说：\BibleQuote{你要按照会幕一切器具的样式竖立会幕，要这样制造。}\BibleRef{出 25:9}不久后又说：\BibleQuote{这样，你要按照在山上向你显示的样式制造。}\BibleRef{出 25:40}\Person{柏拉图}读到这些段落，却没有以适当的洞察力理解所写的内容，便认为理型在感官所察觉的事物之前已有某种独立存在，并常称之为受造之物的样式，因为\Person{梅瑟}关于会幕的记载这样说：\BibleQuote{要按照在山上向你显示的样式制造。}\par

% structure: p0104 source=h2 target=section reason=relative-heading-rank; base=h2

\section{第三十章 荷马对人类起源的知识}

他在关于天、地和人的问题上也同样明显受骗，因为他认为这些都有\OriginalTerm{理型}{ideas}。因为\Person{梅瑟}这样写道：\BibleQuote{起初天主创造了天和地，}接着又补充了这句话：\BibleQuote{地是混沌空虛的，}他以为\BibleQuote{地是}这句话指的是先存的地，因为\Person{梅瑟}说\BibleQuote{地是混沌空虛的；}他以为\Person{梅瑟}所说的\BibleQuote{天主创造了天和地}中的地是我们感官所察觉的地，是天主按照先存的理型所造的。同样，关于被创造的天，他认为那个被创造的天——他也称之为穹苍——是感官所察觉的受造物；而理性所察觉的天则是先知所说的另一个天：\BibleQuote{诸天属于上主，大地却赐给了世人。}关于人也是如此：\Person{梅瑟}先提了人的名字，然后在许多其他创造之后才提到人的形成，说：\BibleQuote{上主天主用地上的尘土造了人。}\BibleRef{创 2:7}因此，他认为最初如此命名的人先于被造的人，而后来用尘土所造的人则是按照先存的理型制造的。\Person{荷马}也从古代神圣的历史中发现了这点——那段历史说\BibleQuote{你本是尘土，仍要归于尘土}\BibleRef{创 3:19}——因而称\Person{赫克托尔}无生命的尸体为哑土。因为他在谴责\Person{阿基琉斯}在死后拖曳\Person{赫克托尔}的尸首时，在某处写道：\par

\Quote{他在哑土上施加侮辱，\newline{}被愤怒蒙蔽了双眼。}\par

又在别处，他让\Person{墨涅拉俄斯}对那些不愿以应有的热忱接受\Person{赫克托尔}单挑挑战的人这样说：\par

\Quote{愿你们都回到土和水去，}\par

在暴怒中把他们分解为原始且最初的泥土构成。\Person{荷马}和\Person{柏拉图}在\Place{埃及}从古代历史中学到这些，并用他们自己的话记录下来。\par

% structure: p0110 source=h2 target=section reason=relative-heading-rank; base=h2

\section{第三十一章 柏拉图熟悉圣经的进一步证据}

因为，若非从阅读先知著作而来，柏拉图怎能知道他所给予我们的信息，即朱庇特在天上驾驭一辆有翼的战车？他从先知关于革鲁宾的以下表述得知此事：\BibleQuote{主的荣耀从殿里出去，停在革鲁宾上；革鲁宾抬起翅膀，轮子也在它们旁边；以色列之天主上主的荣耀在它们之上。}\BibleRef{则 11:22} 采纳这一观念后，夸夸其谈的柏拉图以极大的自信高声喊道：\Quote{伟大的朱庇特确实在天上驾驭他的有翼战车。} 若非从梅瑟和先知那里，他从何处学到这些并如此写作？他又从何处得到他所说的天主存在于一种火实体中的提示？难道不是从列王纪历史第三卷，其中记载：\BibleQuote{主不在风中；风后有地震，但主不在地震中；地震后有火，但主不在火中；火后有轻微细弱的声音。}\BibleRef{列上 19:11} 但虔诚之人必须以深刻而沉思的洞察力从更高意义上来理解这些事。然而柏拉图没有以适当的洞察力注意这些言辞，就说天主存在于一种火实体中。\par

% structure: p0112 source=h2 target=section reason=relative-heading-rank; base=h2

\section{第三十二章 柏拉图关于天上恩赐的学说}

若有人仔细考虑那从天主降临在圣人们身上的恩赐——神圣先知称之为圣神——他将发现，柏拉图在《与梅诺的对话》中以另一个名称宣告了这一点。因为他害怕将天主的恩赐称为\Quote{圣神}，免得他因跟随先知的教导而显得与希腊人为敌，他确实承认这恩赐从天主而来，却不认为宜称其为圣神，而称之为美德。因此，在《与梅诺的对话》中关于回忆，他在就美德提出许多问题之后——无论是可以教导还是不能教导，而必须通过实践获得，还是既不能通过实践也不能通过学习获得，而是人天生的禀赋，还是以某种其他方式而来——他以这些话语作了如下宣告：\Quote{但若我们通过这整个对话正确地进行我们的探究和讨论，美德既非天生的禀赋，也非人可以因教导而领受，而是凭借神圣的命运临到那些临到的人身上。} 我认为，柏拉图从先知那里学习了这些关于圣神的事，便显然将它们转移到他所谓的美德上。因为正如神圣先知所说，同一个神分成为七神，他同样也称同一美德为四美德，千方百计避免提及圣神，却以某种寓言清楚地宣告了先知关于圣神所说的。因为他在《与梅诺的对话》接近结尾处说了这样的话：\Quote{从这一推理，梅诺，美德似乎是凭借神圣的命运临到那些临到的人身上。但关于这一点，我们要清楚地知道，美德以何种方式临到人，当我们作为第一步，将其作为独立探究，着手调查美德本身是什么的时候。} 你看到他是如何只用美德之名来称呼那从上降临的恩赐；然而他认为值得探究，这恩赐被称为美德还是别的什么是否合适，因为他害怕公开称其为圣神，免得他显得是在跟随先知的教导。\par

% structure: p0114 source=h2 target=section reason=relative-heading-rank; base=h2

\section{第三十三章 柏拉图关于时间开端的概念取自梅瑟}

柏拉图从何处得知时间与诸天一同被创造？因为他如此写道：\Quote{时间因此与诸天一同被创造；以便它们一同生成，也一同消解，倘若它们有朝一日消解的话。} 难道他不是从梅瑟的神圣历史中学到这一点的吗？因为他知道时间的创造是由日、月、年获得其原始构成的。由于那与诸天一同被创造的第一天构成了所有时间的开端（因为梅瑟如此写道：\BibleQuote{在起初天主创造了天地}，然后立即加上\BibleQuote{一天成了}，仿佛他要以时间的一部分来指代全部时间），柏拉图称这日为\Quote{时间}，免得他若提到\Quote{日}，就似乎使自己暴露于雅典人的指控，说他完全采用了梅瑟的措辞。他从何处得知他所写的关于诸天消解的事？难道他也不是从神圣先知那里学到这一点，并认为这是他们的学说吗？\par

% structure: p0116 source=h2 target=section reason=relative-heading-rank; base=h2

\section{第三十四章 人为何赋予天主人的形象}

若有人考察肖像之事，并追问最初塑造你们众神的人基于何种理由认为它们具有人的形态，他将发现这也是源自神圣的历史。因为看到梅瑟的历史以天主的口吻说：\BibleQuote{让我们按我们的形象和相似造人}，这些人便以为这意味着人在形态上相似于天主，于是开始这样塑造他们的神，以为他们是在按相似造相似。但是，你们这些希腊人啊，我现在为何要叙述这些事？为使你们知道，不可能从那些人学习真正的宗教，他们甚至在他们赢得异教徒钦佩的那些主题上，也不能写出任何原创的东西，而只是以某种寓言手法在他们自己的著作中提出他们从梅瑟和其他先知那里学到的东西。\par

% structure: p0118 source=h2 target=section reason=relative-heading-rank; base=h2

\section{第三十五章 对希腊人的呼吁}

希腊人啊，时候已经来到，你们既已被世俗历史说服，承认梅瑟和其余先知远比你们中间被视为智者的人更为古老，就应抛弃你们祖先的古老谬误，阅读先知的圣史，从中认识真正的宗教；因为这些经文所呈现给你们的，并非巧言令色的演说，也不是动听而似是而非的言词——这乃是那些想剥夺你们真理之人的特性——而是朴素地使用自然流露的话语和表达，向你们宣告那降临于他们身上的圣神，愿意通过他们，教导那些渴望学习真宗教的人。你们既已抛弃一切虚假的羞耻，以及人类根深蒂固的错误，连同其所有夸张的排场和空洞的喧哗，虽然你们藉此以为拥有了一切优势，但你们应将自身转向那对你们有益的事。因为如果你们现在表现出渴望投向他们错误的完全对立面，你们对你们的父亲也毫无冒犯，因为很有可能他们如今正在阴间哀叹，并以一种太迟的悔改而痛悔；倘若他们有可能从那里向你们显示，在此生结束后所遭遇的事，你们就会知道他们渴望使你们摆脱何等可怕的灾祸。然而，既然在此生中，你们既无法从他们那里得知，也无法从那些在此自称教导那虚假哲学的人那里得知，那么，唯一剩下的事就是：你们必须抛弃你们祖先的错误，阅读圣经作者们的预言，不向他们要求无可挑剔的文辞（因为我们的宗教之事在于行为，而非言语），并从中学习那将赐予你们永生的事。因为那些徒然玷污哲学之名的人，被证实一无所知，连他们自己也难免被迫承认，因为他们不仅彼此意见不合，而且有时用这种方式、有时用那种方式表达自己的观点。\par

% structure: p0120 source=h2 target=section reason=relative-heading-rank; base=h2

\section{第三十六章　哲学家并未拥有真知}

如果在他们中间，\Quote{发现真理}被作为哲学的定义之一，那么那些并未拥有真知识的人，怎配得上哲学之名呢？因为如果苏格拉底——你们智者中最有智慧的人，连你们自己的神谕也为他作证，如你们所说，说\Quote{众人中苏格拉底最有智慧}——他承认自己一无所知，那么那些后来的人，又怎能自称知道天上之事呢？因为苏格拉底说，他之所以被称为有智慧，正是因为别人假装知道他们所不知道的事，而他自己却不畏惧承认自己一无所知。因为他说：\Quote{我似乎因这微小的一点而最有智慧：我不知道的事，我不以为我知道。}不要有人以为苏格拉底是讽刺性地假装无知，因为他常在对话中这样做。因为他在被带往监狱时所说的最后告别之辞，证明他是认真而真实地承认自己的无知：\Quote{但现在该走了，我去赴死，你们去活。我们中谁去往更好的境况，唯有天主知道。}苏格拉底确实在阿勒约帕哥说了这最后一句话，然后便去了监狱，将那些对我们隐藏的事的知识单单归于天主；但那些后来的人，虽然连地上的事都无法理解，却自称明白天上的事，仿佛亲眼见过一样。至少亚里士多德——仿佛比柏拉图更精确地看见了天上的事——宣称天主并不像柏拉图所说的那样存在于火质中（因为这是柏拉图的学说），而是存在于第五元素——空气中。而他虽然要求凭着他语言的优美，人们应当相信他关于这些事的说法，却因为无法发现卡尔基斯欧里普斯海峡的性质而蒙受耻辱和名誉扫地，最终离世。因此，任何有健全判断的人，都不要将这些人的优美言辞置于自己的救恩之上，而应依照那古老的故事，用蜡塞住耳朵，逃避这些塞壬所施予的甜蜜伤害。因为上述之人，以其优美的语言为诱饵，试图引诱许多人离开正确的宗教，效法那胆敢教导最初之人多神论者。我恳求你们，不要被这些人说服，而要阅读圣经作者的预言。如果任何怠惰或古老的世袭迷信阻止你们阅读圣人们的预言——通过它们你们可以受教认识唯一的天主，这是真宗教的首要信条——那么，至少要相信那位最初教导你们多神论，后来却又宁愿唱出有益而必要的悔改之歌的人，即俄耳甫斯，他所说的正是我前文所引的话；也要相信那些就唯一的天主写了同样内容的其他作者。因为这是天主眷顾为你们所做的工作：他们虽不情愿，却为主张了先知们关于唯一的天主所说的话是真实的，为的是，在多神论的教义被众人拒绝之后，你们有机会认识真理。\par

% structure: p0122 source=h2 target=section reason=relative-heading-rank; base=h2

\section{第三十七章　论西彼拉}

你们可以部分地轻易从古老的\Person{西比尔}学到真宗教，她通过某种强大的默感，通过她的神谕预言教导你们一些似乎与先知们的教导十分相近的真理。他们说，她是巴比伦血统，是\Person{贝罗苏斯}的女儿，\Person{贝罗苏斯}曾写了\WorkTitle{迦勒底史}；她（不知如何）越过到了\Place{坎帕尼亚}地区，在那里一个叫\Place{库迈}的城市发表她的神谕，\Place{库迈}距\Place{巴亚}六英里，\Place{坎帕尼亚}的温泉就在那里。我们曾在那城里，也看见了一个地方，有人指给我们看一个由一整块石头凿出的非常大的大殿；一个巨大的建筑，值得一切赞叹。从父辈那里听到这作为国家传统一部分的人告诉我们，她曾在这里发布她的神谕。在大殿中央，他们指给我们看三个由一块石头凿出的容器，他们说，当里面装满水时，她就在那里洗澡，然后重新穿上袍子，隐退到大殿最里面的房间，那房间仍然是一块石头的一部分；她坐在房间中央的高讲坛和宝座上，这样宣布她的神谕。\Person{西比尔}作为女先知曾被许多其他作家提及，也被\Person{柏拉图}在他的\WorkTitle{斐德若}中提及。\Person{柏拉图}似乎在她阅读她的预言时，将先知视为神圣默感的人。因为他看见她很久以前预言的已经实现；因此他在\WorkTitle{美诺篇}中以如下言辞表达了对先知的惊奇和钦佩：\Quote{我们现在称为先知的人，我们应当正当地称他们为神圣的。而且我们至少会说，他们是神圣的，并且被天主的灵感和拥有提升到先知的狂喜，当他们正确地谈到许多重要的事情，却对所说的一无所知，}——清楚而明显地指的是\Person{西比尔}的预言。因为不同于诗人，诗人在他们的诗篇写就后有权力修改和润色，特别是增加他们诗句的准确性，而她在灵感降临时确实充满了预言，但灵感一停止，她所说的一切的回忆也随之停止。这确实正是\Person{西比尔}的一些诗句（而非全部）的格律得以保存的原因。因为当我们自己在那个城市时，从我们的\Emph{导游}（他带领我们看了她过去预言的地方）得知，有一个黄铜做的箱子，他们说她的遗骸保存在里面。除了他们从父辈那里听到的一切之外，他们还告诉我们，当时记录她预言的人，是未受教育的人，常常在格律的准确性上完全出错；他们说，这就是一些诗句缺乏格律的原因，因为女先知在附身和灵感停止后对所说过的话已无记忆，而记录者因缺乏教育未能准确地记录格律。因此，显然\Person{柏拉图}在说关于先知的话时，心中想着\Person{西比尔}的预言，因为他曾说：当他们正确地谈到许多重要的事情，却对所说的一无所知。\par

% structure: p0124 source=h2 target=section reason=relative-heading-rank; base=h2

\section{第三十八章 结语呼吁}

但是，希腊人啊，既然真宗教的事不在于诗歌的韵律数目，也不在于你们中间备受推崇的文化，你们今后应较少致力于格律和语言的准确性，而不加争辩地留意\Person{西比尔}的话，要认出她将以清晰显明的方式预言我们的救主耶稣基督的来临，将赐予你们多大的益处。祂作为天主的圣言，在能力上与祂不可分离，取了按天主的肖像和模样所造的人，恢复了我们古老祖先的宗教知识，他们之后的人因嫉妒魔鬼的蛊惑计谋而离弃了这知识，转而敬拜那些不是神的神。如果你们仍然犹豫，因人的形成之事而被阻碍相信，就请相信那些你们迄今认为应当留意的人，要知道你们自己的神谕，当有人请求它向全能天主唱一篇赞美诗时，它在颂歌中间这样说了：\Quote{谁造了最初的人，称他为亚当。}这首赞美诗被许多我们所认识的人保存着，为要说服那些不愿相信众人所见证的真理的人。因此，希腊人啊，如果你们不把那些关于不是神之神的虚假幻想看得比你们自己的救恩更重，那么，如我所说，请相信那位最古老且备受尊崇的\Person{西比尔}，她的书在全世界被保存；她通过某种强大的默感，既在她的神谕话语中教导我们关于那些被称为神但实际不存在的事；也清楚而明显地预言了我们的救主耶稣基督被预言的来临，以及祂所要做的所有事情。因为这些知识将构成你们学习圣经作者预言的必要预备训练。如果有人以为自己从那些你们称为哲学家中最古老的人那里学到了关于天主的道理，就让他听\Person{阿蒙}和\Person{赫尔墨斯}的话：\Person{阿蒙}在他的论天主的话语中称祂为全然隐藏；\Person{赫尔墨斯}则清楚而明确地说，\Quote{理解天主是困难的，甚至连能理解祂的人也不可能向别人宣告祂。}因此，从各方面来看，必须看到，除了通过那些由神圣默感教导我们的先知以外，根本不可能学到任何关于天主和真正宗教的事。\par

\SourceRecord{Translated by Marcus Dods.；Ante-Nicene Fathers；Vol. 1.；Edited by Alexander Roberts, James Donaldson, and A. Cleveland Coxe.；Buffalo, NY: Christian Literature Publishing Co.,；1885.；<http://www.newadvent.org/fathers/0129.htm>.}
